%============================================================
%  Bd. 22 - Halbgruppen mit Nullelement
%============================================================

\documentclass[main.tex]{subfiles}

\ifSubfilesClassLoaded{
  \BandSetupStandalone{B22}
}{
}

\title{Bd. 22 - Halbgruppen mit Nullelement}
\author{Martin Kunze}
\date{}
\setcounter{file}{22}

\hypersetup{
  pdftitle={Bd. 22 - Halbgruppen mit Nullelement},
  pdfauthor={Martin Kunze}
}

\begin{document}

\title{Bd. 22 - Halbgruppen mit Nullelement}
\author{Martin Kunze}
\date{}
\hypersetup{pageanchor=false}
\maketitle
\hypersetup{pageanchor=true}
\setcounter{tocdepth}{3}
\tableofcontents
\setcounter{file}{22}
\renewcommand{\FormulaBandID}{22}

\chapter{Einführung}

Dieser Band untersucht Halbgruppen mit einem absorbierenden Nullelement.

\chapter{Halbgruppen mit Nullelement}

\section{Definition}

Eine Halbgruppe kann ein absorbierendes Element besitzen.  Wir nennen es
\emph{Nullelement}; der Ausdruck \emph{Nullhalbgruppe} bleibt dagegen der
spezielleren Situation vorbehalten, in der jedes Produkt gleich dem
Nullelement ist.

\FormulaDefDeltaK[Halbgruppe mit Nullelement]{%
  \SemigroupWithZeroAlg{A}{\star}{0}%
}{SemigroupWithZeroDef}{%
  \DeltaRow{Mengen}{A\dsep 0\dsep x}
  \DeltaRow{\textbf{Axiome}}{}
  \DeltaPrem{Halbgruppe}{\SemigroupAlg{A}{\star}}
  \DeltaPrem{Nullelement}{0\in A}
  \DeltaRow{Linksabsorption}{x\in A\vdash 0\star x=0}
  \DeltaRow{Rechtsabsorption}{x\in A\vdash x\star0=0}
  \DeltaRow{\textbf{Neues Symbol}}{}
  \DeltaPrem{Halbgruppe mit Nullelement}{%
    \SemigroupWithZeroAlg{A}{\star}{0}}
}

\FormulaAxiomDeltaK[Zugriff auf die Halbgruppe]{%
  \SemigroupWithZeroAlg{A}{\star}{0}
  \vdash\SemigroupAlg{A}{\star}%
}{SemigroupWithZeroBaseAxiom}{%
  \DeltaRow{Mengen}{A\dsep 0}
}

\FormulaAxiomDeltaK[Zugriff auf das Nullelement]{%
  \SemigroupWithZeroAlg{A}{\star}{0}\vdash 0\in A%
}{SemigroupWithZeroCarrierAxiom}{%
  \DeltaRow{Mengen}{A\dsep 0}
}

\FormulaAxiomDeltaK[Zugriff auf die Linksabsorption]{%
  \SemigroupWithZeroAlg{A}{\star}{0}\dsep x\in A
  \vdash 0\star x=0%
}{SemigroupWithZeroLeftAbsorptionAxiom}{%
  \DeltaRow{Mengen}{A\dsep 0\dsep x}
}

\FormulaAxiomDeltaK[Zugriff auf die Rechtsabsorption]{%
  \SemigroupWithZeroAlg{A}{\star}{0}\dsep x\in A
  \vdash x\star0=0%
}{SemigroupWithZeroRightAbsorptionAxiom}{%
  \DeltaRow{Mengen}{A\dsep 0\dsep x}
}

\FormulaThmDeltaK[Eindeutigkeit des Nullelements]{%
  \SemigroupWithZeroAlg{A}{\star}{0}\dsep
  \SemigroupWithZeroAlg{A}{\star}{z}
  \vdash 0=z%
}{SemigroupZeroUnique}{%
  \DeltaRow{Mengen}{A\dsep 0\dsep z}
  \DeltaRow{Funktionen}{\star}
}
\begin{tabproof}
  \proofstep{1}{\SemigroupWithZeroAlg{A}{\star}{0}}{\rA}
  \proofstep{2}{\SemigroupWithZeroAlg{A}{\star}{z}}{\rA}
  \proofstep{1}{0\in A}{%
    \FormulaRefAuto{SemigroupWithZeroCarrierAxiom}[ax]{1}}
  \proofstep{2}{z\in A}{%
    \FormulaRefAuto{SemigroupWithZeroCarrierAxiom}[ax]{2}}
  \proofstep{1,2}{0\star z=0}{%
    \FormulaRefAuto{SemigroupWithZeroLeftAbsorptionAxiom}[ax]{1,4}}
  \proofstep{1,2}{0\star z=z}{%
    \FormulaRefAuto{SemigroupWithZeroRightAbsorptionAxiom}[ax]{2,3}}
  \proofstep{1,2}{0=z}{%
    \FormulaRefAuto{a = b,\, a = c \vdash b = c}[thm]{5,6}}
\end{tabproof}

\section{Beispiele}

\FormulaThmDeltaK[Die Multiplikation der natürlichen Zahlen besitzt ein
Nullelement]{%
  \SemigroupWithZeroAlg{\mathbb{N}}{\cdot}{0}%
}{NaturalMultiplicationSemigroupWithZero}{%
  \DeltaRow{Mengen}{\mathbb{N}\dsep 0\dsep a}
}
\begin{tabproof}
  \proofstep{}{\SemigroupAlg{\mathbb{N}}{\cdot}}{%
    \FormulaRefAuto{NaturalMultiplicationSemigroup}[thm]}
  \proofstep{}{0\in\mathbb{N}}{%
    \FormulaRefAuto{PeanoZero}[ax]}
  \proofstep{3}{a\in\mathbb{N}}{\rA}
  \proofstep{3}{0\cdot a=0}{%
    \FormulaRefAuto{PeanoMulLeftZero}[thm]{3}}
  \proofstep{3}{a\cdot0=0}{%
    \FormulaRefAuto{PeanoMulRightZero}[thm]{3}}
  \proofstep{}{\SemigroupWithZeroAlg{\mathbb{N}}{\cdot}{0}}{%
    \FormulaRefAuto{SemigroupWithZeroDef}[def]{1,2,4,5}}
\end{tabproof}

\FormulaThmDeltaK[Eine konstante Operation bildet eine Halbgruppe mit
Nullelement]{%
  \begin{aligned}[t]
    &p\in A\dsep
      \forall x,y\in A\,(x\star_0y=p)\vdash{}\\[-2pt]
    &\SemigroupWithZeroAlg{A}{\star_0}{p}
  \end{aligned}
}{ConstantSemigroupWithZero}{%
  \DeltaRow{Mengen}{A\dsep p\dsep x\dsep y\dsep z}
  \DeltaRow{Funktionen}{\star_0}
}
\begin{tabproof}
  \proofstep{1}{p\in A}{\rA}
  \proofstep{2}{\forall x,y\in A\,(x\star_0y=p)}{\rA}
  \proofstep{3}{x\in A}{\rA}
  \proofstep{4}{y\in A}{\rA}
  \proofstep{5}{z\in A}{\rA}
  \proofstep{2,3,4}{x\star_0y=p}{\rRE{3,\rRE{4,2}}}
  \proofstep{1,2,3,4}{x\star_0y\in A}{\rIE{6,1}}
  \proofstep{1,2,3,4,5}{(x\star_0y)\star_0z=p}{%
    \rRE{7,\rRE{5,2}}}
  \proofstep{2,4,5}{y\star_0z=p}{\rRE{4,\rRE{5,2}}}
  \proofstep{1,2,4,5}{y\star_0z\in A}{\rIE{9,1}}
  \proofstep{1,2,3,4,5}{x\star_0(y\star_0z)=p}{%
    \rRE{3,\rRE{10,2}}}
  \proofstep{1,2,3,4,5}{p=x\star_0(y\star_0z)}{%
    \FormulaRefAuto{a=b\vdash b=a}[thm]{11}}
  \proofstep{1,2,3,4,5}{%
    (x\star_0y)\star_0z=x\star_0(y\star_0z)}{\rChain{8,12}}
  \proofstep{1,2}{\SemigroupAlg{A}{\star_0}}{%
    \FormulaRefAuto{SemigroupDef}[def]{7,13}}
  \proofstep{1,2,3}{p\star_0x=p}{\rRE{1,\rRE{3,2}}}
  \proofstep{1,2,3}{x\star_0p=p}{\rRE{3,\rRE{1,2}}}
  \proofstep{1,2}{\SemigroupWithZeroAlg{A}{\star_0}{p}}{%
    \FormulaRefAuto{SemigroupWithZeroDef}[def]{14,1,15,16}}
\end{tabproof}

\section{Annihilatoren}

Annihilatoren werden für Halbgruppen mit Nullelement und relativ zu einer
Teilmenge definiert.  Links und rechts müssen getrennt werden, solange keine
Kommutativität vorausgesetzt ist.

\FormulaDefDeltaK[Linksannihilator einer Teilmenge]{%
  \begin{aligned}[t]
    &\SemigroupWithZeroAlg{A}{\star}{0}\dsep X\subseteq A\vdash{}\\[-2pt]
    &\SemigroupLeftAnnihilator{A,\star,0}{X}
      \coloneqq
      \{a\in A\mid\forall x\in X\,(a\star x=0)\}
  \end{aligned}%
}{SemigroupLeftAnnihilatorDef}{%
  \DeltaRow{Mengen}{A\dsep X\dsep 0\dsep a\dsep x}
  \DeltaRow{Funktionen}{\star}
  \DeltaRow{\textbf{Neues Symbol}}{}
  \DeltaPrem{Linksannihilator}{%
    \SemigroupLeftAnnihilator{A,\star,0}{X}}
}

\FormulaDefDeltaK[Rechtsannihilator einer Teilmenge]{%
  \begin{aligned}[t]
    &\SemigroupWithZeroAlg{A}{\star}{0}\dsep X\subseteq A\vdash{}\\[-2pt]
    &\SemigroupRightAnnihilator{A,\star,0}{X}
      \coloneqq
      \{a\in A\mid\forall x\in X\,(x\star a=0)\}
  \end{aligned}%
}{SemigroupRightAnnihilatorDef}{%
  \DeltaRow{Mengen}{A\dsep X\dsep 0\dsep a\dsep x}
  \DeltaRow{Funktionen}{\star}
  \DeltaRow{\textbf{Neues Symbol}}{}
  \DeltaPrem{Rechtsannihilator}{%
    \SemigroupRightAnnihilator{A,\star,0}{X}}
}

\newpage

\FormulaDefDeltaK[Zweiseitiger Annihilator einer Teilmenge]{%
  \begin{aligned}[t]
    &\SemigroupWithZeroAlg{A}{\star}{0}\dsep X\subseteq A\vdash{}\\[-2pt]
    &\SemigroupAnnihilator{A,\star,0}{X}
      \coloneqq
      \SemigroupLeftAnnihilator{A,\star,0}{X}
      \cap\SemigroupRightAnnihilator{A,\star,0}{X}
  \end{aligned}%
}{SemigroupAnnihilatorDef}{%
  \DeltaRow{Mengen}{A\dsep X\dsep 0}
  \DeltaRow{Funktionen}{\star}
  \DeltaRow{\textbf{Neues Symbol}}{}
  \DeltaPrem{Zweiseitiger Annihilator}{%
    \SemigroupAnnihilator{A,\star,0}{X}}
}

Damit gilt ausgeschrieben
\[
  \SemigroupAnnihilator{A,\star,0}{X}
  =\{a\in A\mid
      \forall x\in X\,(a\star x=0\land x\star a=0)\}.
\]
Für \(X=A\) heißt
\(\SemigroupAnnihilator{A,\star,0}{A}\) der
\emph{Totalannihilator}.  Bei Kommutativität stimmen Links-, Rechts- und
zweiseitiger Annihilator überein.  Insbesondere gehört das Nullelement zu
jedem dieser Annihilatoren.

\section{Morphismen}

\FormulaDefDeltaK[Homomorphismus von Halbgruppen mit Nullelement]{%
  \SemigroupWithZeroHom{f}{A,\star,0_A}{B,\diamond,0_B}%
}{SemigroupWithZeroHomDef}{%
  \DeltaRow{Mengen}{A\dsep B\dsep0_A\dsep0_B}
  \DeltaRow{Funktionen}{f\dsep\star\dsep\diamond}
  \DeltaRow{\textbf{Axiome}}{}
  \DeltaPrem{Quellstruktur}{\SemigroupWithZeroAlg{A}{\star}{0_A}}
  \DeltaPrem{Zielstruktur}{\SemigroupWithZeroAlg{B}{\diamond}{0_B}}
  \DeltaPrem{Halbgruppenhomomorphismus}{%
    \SemigroupHom{f}{A,\star}{B,\diamond}}
  \DeltaRow{Nullerhaltung}{f(0_A)=0_B}
  \DeltaRow{\textbf{Neues Symbol}}{}
  \DeltaPrem{Homomorphismus}{%
    \SemigroupWithZeroHom{f}{A,\star,0_A}{B,\diamond,0_B}}
}

\FormulaAxiomDeltaK[Quellstruktur]{%
  \SemigroupWithZeroHom{f}{A,\star,0_A}{B,\diamond,0_B}
  \vdash\SemigroupWithZeroAlg{A}{\star}{0_A}%
}{SemigroupWithZeroHomSourceAxiom}{%
  \DeltaRow{Mengen}{A\dsep B\dsep0_A\dsep0_B}
  \DeltaRow{Funktionen}{f\dsep\star\dsep\diamond}
}

\FormulaAxiomDeltaK[Zielstruktur]{%
  \SemigroupWithZeroHom{f}{A,\star,0_A}{B,\diamond,0_B}
  \vdash\SemigroupWithZeroAlg{B}{\diamond}{0_B}%
}{SemigroupWithZeroHomTargetAxiom}{%
  \DeltaRow{Mengen}{A\dsep B\dsep0_A\dsep0_B}
  \DeltaRow{Funktionen}{f\dsep\star\dsep\diamond}
}

\FormulaAxiomDeltaK[Zugriff auf den Halbgruppenhomomorphismus]{%
  \SemigroupWithZeroHom{f}{A,\star,0_A}{B,\diamond,0_B}
  \vdash\SemigroupHom{f}{A,\star}{B,\diamond}%
}{SemigroupWithZeroHomBaseHomAxiom}{%
  \DeltaRow{Mengen}{A\dsep B\dsep0_A\dsep0_B}
  \DeltaRow{Funktionen}{f\dsep\star\dsep\diamond}
}

\FormulaAxiomDeltaK[Zugriff auf die Nullerhaltung]{%
  \SemigroupWithZeroHom{f}{A,\star,0_A}{B,\diamond,0_B}
  \vdash f(0_A)=0_B%
}{SemigroupWithZeroHomZeroAxiom}{%
  \DeltaRow{Mengen}{A\dsep B\dsep0_A\dsep0_B}
  \DeltaRow{Funktionen}{f\dsep\star\dsep\diamond}
}

\FormulaDefDeltaK[Isomorphismus von Halbgruppen mit Nullelement]{%
  \SemigroupWithZeroIso{f}{A,\star,0_A}{B,\diamond,0_B}%
}{SemigroupWithZeroIsoDef}{%
  \DeltaRow{Mengen}{A\dsep B\dsep0_A\dsep0_B}
  \DeltaRow{Funktionen}{f\dsep\star\dsep\diamond}
  \DeltaRow{\textbf{Axiome}}{}
  \DeltaPrem{Homomorphismus}{%
    \SemigroupWithZeroHom{f}{A,\star,0_A}{B,\diamond,0_B}}
  \DeltaPrem{Bijektivität}{f\colon A\bij B}
  \DeltaRow{\textbf{Neues Symbol}}{}
  \DeltaPrem{Isomorphismus}{%
    \SemigroupWithZeroIso{f}{A,\star,0_A}{B,\diamond,0_B}}
}

\FormulaAxiomDeltaK[Homomorphismuszugriff]{%
  \SemigroupWithZeroIso{f}{A,\star,0_A}{B,\diamond,0_B}
  \vdash\SemigroupWithZeroHom{f}{A,\star,0_A}{B,\diamond,0_B}%
}{SemigroupWithZeroIsoHomAxiom}{%
  \DeltaRow{Mengen}{A\dsep B\dsep0_A\dsep0_B}
  \DeltaRow{Funktionen}{f\dsep\star\dsep\diamond}
}

\FormulaAxiomDeltaK[Bijektivitätszugriff]{%
  \SemigroupWithZeroIso{f}{A,\star,0_A}{B,\diamond,0_B}
  \vdash f\colon A\bij B%
}{SemigroupWithZeroIsoBijectiveAxiom}{%
  \DeltaRow{Mengen}{A\dsep B\dsep0_A\dsep0_B}
  \DeltaRow{Funktionen}{f\dsep\star\dsep\diamond}
}

\end{document}
